\documentclass{article}
\usepackage[T1]{fontenc}
\usepackage{lmodern}
\usepackage{graphicx}
\usepackage{amsmath,amssymb}
\usepackage[a4paper,margin=2cm]{geometry}
\usepackage[absolute,overlay]{textpos}
\setlength{\TPHorizModule}{1mm}
\setlength{\TPVertModule}{1mm}
\usepackage[indonesian]{babel}
\usepackage{hyperref}
\setlength{\emergencystretch}{2em}
% unit: Halaman 52

\begin{document}

% === Halaman 52 (python/native) ===

52

• Ukuran pesan yang akan disembunyikan bergantung pada ukuran cover-object.

• Misalkan pada citra grayscale (1 byte/pixel) 256 x 256 pixel :

 - jumlah pixel = jumlah byte = 256 x 256 = 65536

 - setiap byte dapat menyembunyikan 1 bit pesan di LSB-nya

 - jadi ukuran maksimal pesan = 65536 bit = 8192 byte = 8 KB

• Pada citra berwarna 24-bit berukuran 256  256 pixel:

 - jumlah pixel 256 x 256 = 65536

 - setiap pixel = 3 byte, berarti ada 65536  3 = 196608 byte. 

 - setiap byte dapat menyembunyikan 1 bit pesan

 - jadi ukuran maksimal pesan = 196608 bit = 24576 byte = 24KB

Menghitung Ukuran Pesan yang dapat Disembunyikan

\end{document}
